\section{Discusión}

El análisis comparativo de los resultados obtenidos en esta simulación frente a la literatura científica especializada reveló importantes coincidencias y divergencias metodológicas que enriquecen la comprensión de la seguridad en redes académicas.

En primer lugar, la simulación de propagación lateral en el Escenario A mostró una curva de infección de carácter exponencial que condujo al colapso total de la red. Este comportamiento se contrastó con lo establecido por Chernikova et al. \cite{chernikova2023modeling}, quienes sostienen que los modelos deterministas tradicionales como el SIR (Susceptible-Infectado-Recuperado) no logran capturar con precisión la latencia del malware moderno en redes complejas, por lo cual proponen el modelo SIIDR (Susceptible-Infectado-Infectado Latente-Recuperado) para introducir estados durmientes. En la simulación desarrollada en GAMA, se implementó un modelo de transición reactiva directa sin periodo de latencia. Si bien esta simplificación incrementó la agresividad inicial del ataque, facilitó la medición del peor escenario de propagación lateral.

En segundo lugar, se observó que la topología física y lógica de la red universitaria (estructurada mediante laboratorios conectados en cascada hacia un switch central) actuó como un vector facilitador para que la infección se dispersase de forma secuencial dentro de cada sala antes de saltar a otros laboratorios. Este hallazgo valida los análisis estructurales presentados en IEEE Access \cite{ieee2019spreading}, donde se demuestra que la matriz de adyacencia de una red dicta el umbral de propagación del ransomware y que los nodos concentradores o \textit{hubs} (como el switch principal en este modelo) representan puntos críticos de falla sistémica que aceleran de forma exponencial la tasa de transmisión global si no disponen de mecanismos de segmentación lógica.

Finalmente, la efectividad del protocolo de contingencia activado en el Escenario B, el cual detuvo la infección al alcanzar el 30\% de equipos comprometidos, coincidió con las simulaciones de control óptimo de epidemias de red de IEEE Access \cite{ieee2025optimal}. En dicho estudio, los autores demuestran que las intervenciones de parcheo dinámico y cuarentena selectiva reducen significativamente el costo de recuperación y el daño operativo si se aplican en las etapas tempranas de un brote. Esto dio sustento a los resultados de nuestro experimento, donde la acción reactiva temprana preservó el 70\% de la red LAN. Además, esta resiliencia basada en múltiples capas de seguridad concuerda con el modelo de riesgo cibernético de Padilla et al. \cite{padilla2025serdux}, quienes señalan que la efectividad de una defensa no depende únicamente del perímetro (firewall), sino de la capacidad de respuesta y aislamiento local en las terminales una vez que la barrera externa es superada.
